%%%%%%%%%%%%%%%%%%%%%%%%%%%%%%%%%%%%%%%%%%%%%%%%%%%%%%%%%%%%%%%%%%%%%%
% Lớp 8 | Chương IV. Định lí Thalès | Định lí Thalès trong tam giác
% Dạng 8H4?15-3: Định lí Thalès đảo: chứng minh song song
%%%%%%%%%%%%%%%%%%%%%%%%%%%%%%%%%%%%%%%%%%%%%%%%%%%%%%%%%%%%%%%%%%%%%%


% nguồn: dan_tu_AI (nhập 16/11/2026)
\begin{bt}%[8H4H15-3]%[ma:c4878a84]
	Cho tam giác $ABC$. Trên cạnh $AB$ lấy điểm $D$, trên cạnh $AC$ lấy điểm $E$ sao cho $AD=3\text{ cm}$, $DB=6\text{ cm}$, $AE=4\text{ cm}$, $EC=8\text{ cm}$. Chứng minh $DE\parallel BC$.
	\loigiai{
		Ta có $\dfrac{AD}{DB}=\dfrac36=\dfrac12$ và $\dfrac{AE}{EC}=\dfrac48=\dfrac12$.\\
		Suy ra $\dfrac{AD}{DB}=\dfrac{AE}{EC}$.\\
		Theo định lí Thalès đảo, $DE\parallel BC$.
	}
\end{bt}

% nguồn: dan_tu_AI (nhập 16/11/2026)
\begin{bt}%[8H4H15-3]%[ma:a19e87af]
	Cho tam giác $ABC$ có $AB=12\text{ cm}$, $AC=15\text{ cm}$. Trên cạnh $AB$ lấy điểm $D$ sao cho $AD=4\text{ cm}$, trên cạnh $AC$ lấy điểm $E$ sao cho $AE=5\text{ cm}$. Chứng minh $DE\parallel BC$.
	\loigiai{
		Ta có $\dfrac{AD}{AB}=\dfrac{4}{12}=\dfrac13$ và $\dfrac{AE}{AC}=\dfrac{5}{15}=\dfrac13$.\\
		Suy ra $\dfrac{AD}{AB}=\dfrac{AE}{AC}$.\\
		Theo định lí Thalès đảo, $DE\parallel BC$.
	}
\end{bt}

% nguồn: dan_tu_AI (nhập 16/11/2026)
\begin{bt}%[8H4H15-3]%[ma:5acd90cc]
	Cho tam giác $ABC$, điểm $M$ thuộc cạnh $AB$, điểm $N$ thuộc cạnh $AC$ với $AM=2\text{ cm}$, $MB=3\text{ cm}$, $AN=3\text{ cm}$, $NC=4\text{ cm}$. Đường thẳng $MN$ có song song với $BC$ không? Vì sao?
	\loigiai{
		Ta có $\dfrac{AM}{MB}=\dfrac23$ và $\dfrac{AN}{NC}=\dfrac34$.\\
		Vì $\dfrac23\ne\dfrac34$ nên $\dfrac{AM}{MB}\ne\dfrac{AN}{NC}$.\\
		Nếu $MN\parallel BC$ thì theo định lí Thalès phải có $\dfrac{AM}{MB}=\dfrac{AN}{NC}$, mâu thuẫn.\\
		Vậy $MN$ không song song với $BC$.
	}
\end{bt}

% nguồn: dan_tu_AI (nhập 16/11/2026)
\begin{bt}%[8H4H15-3]%[ma:6699e5ec]
	Cho tam giác $ABC$ có $AB=9\text{ cm}$. Trên cạnh $AB$ lấy điểm $M$ sao cho $AM=6\text{ cm}$; trên cạnh $AC$ lấy điểm $N$ sao cho $AN=8\text{ cm}$, $NC=4\text{ cm}$. Chứng minh $MN\parallel BC$.
	\loigiai{
		$AC=AN+NC=8+4=12\text{ (cm)}$.\\
		Ta có $\dfrac{AM}{AB}=\dfrac69=\dfrac23$ và $\dfrac{AN}{AC}=\dfrac{8}{12}=\dfrac23$.\\
		Suy ra $\dfrac{AM}{AB}=\dfrac{AN}{AC}$. Theo định lí Thalès đảo, $MN\parallel BC$.
	}
\end{bt}

% nguồn: dan_tu_AI (nhập 16/11/2026)
\begin{bt}%[8H4H15-3]%[ma:6da49632]
	Cho tam giác $ABC$. Lấy điểm $D$ trên cạnh $AB$, điểm $E$ trên cạnh $AC$, điểm $F$ trên cạnh $BC$ sao cho $\dfrac{AD}{DB}=\dfrac{AE}{EC}=\dfrac12$ và $\dfrac{BF}{FC}=2$.\\
	a) Chứng minh $DE\parallel BC$ và $DF\parallel AC$.\\
	b) Tứ giác $DECF$ là hình gì? Vì sao?
	\loigiai{
		a) Trong tam giác $ABC$: $\dfrac{AD}{DB}=\dfrac{AE}{EC}$ nên $DE\parallel BC$ (định lí Thalès đảo).\\
		Từ $\dfrac{AD}{DB}=\dfrac12$ suy ra $\dfrac{BD}{BA}=\dfrac23$. Từ $\dfrac{BF}{FC}=2$ suy ra $\dfrac{BF}{BC}=\dfrac23$.\\
		Trong tam giác $BAC$: $\dfrac{BD}{BA}=\dfrac{BF}{BC}$ nên $DF\parallel AC$ (định lí Thalès đảo).\\
		b) Tứ giác $DECF$ có $DE\parallel FC$ (vì $DE\parallel BC$) và $DF\parallel EC$ (vì $DF\parallel AC$).\\
		Vậy $DECF$ là hình bình hành.
	}
\end{bt}

% nguồn: dan_tu_AI (nhập 16/11/2026)
\begin{bt}%[8H4V15-3]%[ma:c87684bd]
	Cho tam giác $ABC$ có đường trung tuyến $AM$. Trên cạnh $AB$ lấy điểm $D$, trên cạnh $AC$ lấy điểm $E$ sao cho $\dfrac{AD}{AB}=\dfrac{AE}{AC}$. Đoạn $AM$ cắt $DE$ tại $I$.\\
	a) Chứng minh $DE\parallel BC$.\\
	b) Chứng minh $I$ là trung điểm của $DE$.
	\loigiai{
		a) Vì $\dfrac{AD}{AB}=\dfrac{AE}{AC}$ nên theo định lí Thalès đảo, $DE\parallel BC$.\\
		b) Tam giác $ABM$ có $DI\parallel BM$ nên theo hệ quả của định lí Thalès: $\dfrac{DI}{BM}=\dfrac{AI}{AM}$.\\
		Tam giác $ACM$ có $IE\parallel MC$ nên theo hệ quả của định lí Thalès: $\dfrac{IE}{MC}=\dfrac{AI}{AM}$.\\
		Suy ra $\dfrac{DI}{BM}=\dfrac{IE}{MC}$. Mà $BM=MC$ ($M$ là trung điểm $BC$) nên $DI=IE$.\\
		Vậy $I$ là trung điểm của $DE$.
	}
\end{bt}

% nguồn: dan_tu_AI (nhập 18/11/2026)
\begin{ex}%[8H4N15-3]%[ma:75bdd22e]
	Cho tam giác $ABC$, các điểm $D, E$ lần lượt nằm trên $AB, AC$ sao cho $\dfrac{AD}{AB} = \dfrac{AE}{AC}$. Khẳng định nào sau đây đúng?
	\choice
	{$DE \perp BC$}
	{$DE$ cắt $BC$}
	{\True $DE \parallel BC$}
	{$AD = AE$}
	\loigiai{
		Theo định lí Thalès đảo, $\dfrac{AD}{AB} = \dfrac{AE}{AC} \Rightarrow DE \parallel BC$.
	}
\end{ex}

% nguồn: dan_tu_AI (nhập 18/11/2026)
\begin{ex}%[8H4V15-3]%[ma:e5b4511b]
	Cho tam giác $ABC$ có $AB=12\text{ cm}$, $AC=18\text{ cm}$. Trên $AB, AC$ lần lượt lấy $M, N$ sao cho $AM=4\text{ cm}$, $AN=6\text{ cm}$. Gọi $I$ là giao điểm của trung tuyến $AL$ và $MN$. Khẳng định nào sau đây là sai?
	\choice
	{$MN \parallel BC$}
	{\True $MI = 2IL$}
	{$\dfrac{AM}{MB} = \dfrac{AN}{NC}$}
	{$MI = IN$}
	\loigiai{
		Ta có $\dfrac{AM}{AB} = \dfrac{4}{12} = \dfrac{1}{3}$ và $\dfrac{AN}{AC} = \dfrac{6}{18} = \dfrac{1}{3} \Rightarrow MN \parallel BC$.
		Trong $\triangle ABL$, do $MI \parallel BL \Rightarrow \dfrac{MI}{BL} = \dfrac{AI}{AL}$.
		Tương tự $\dfrac{IN}{LC} = \dfrac{AI}{AL}$. Mà $BL = LC \Rightarrow MI = IN$.
		Khẳng định $MI = 2IL$ không có cơ sở đúng.
	}
\end{ex}

% nguồn: dan_tu_AI (nhập 18/11/2026)
\begin{ex}%[8H4C15-3]%[ma:47336e1c]
	Cho tam giác $ABC$, điểm $D$ thuộc tia đối của tia $CB$ sao cho $CD = \dfrac{1}{2}BC$. Điểm $E$ thuộc cạnh $AC$ sao cho $CE = \dfrac{1}{3}AC$. Gọi $I$ là giao điểm của $AB$ và $DE$. Tỉ số $\dfrac{AI}{IB}$ là
	\choice
	{\True $1$}
	{$\dfrac{1}{2}$}
	{$2$}
	{$\dfrac{2}{3}$}
	\loigiai{
		Kẻ $CF \parallel AB$ ($F \in DE$).
		Theo Thalès: $\dfrac{CF}{IB} = \dfrac{CD}{BD} = \dfrac{\frac{1}{2}BC}{\frac{3}{2}BC} = \dfrac{1}{3}$.
		Và $\dfrac{CF}{AI} = \dfrac{CE}{AE} = \dfrac{\frac{1}{3}AC}{\frac{2}{3}AC} = \dfrac{1}{2} \Rightarrow CF = \dfrac{1}{2}AI$.
		Suy ra $\dfrac{1}{3}IB = \dfrac{1}{2}AI \Rightarrow \dfrac{AI}{IB} = \dfrac{2}{3}$.
		(Khoan, kiểm tra lại: $\dfrac{CF}{IB} = 1/3 \Rightarrow IB = 3CF$. $\dfrac{CF}{AI} = 1/2 \Rightarrow AI = 2CF$. Vậy $\dfrac{AI}{IB} = \dfrac{2}{3}$.
		Phải đổi phương án đúng thành 2/3. Do yêu cầu đảo rải đáp án, ta viết lại lời giải chuẩn).
		% KIEM_TRA: Sửa lại tỉ lệ cho đúng.
		Kẻ $CK \parallel AB$ ($K \in DE$). $\dfrac{CK}{IB} = \dfrac{CD}{BD} = \dfrac{1/2}{3/2} = \dfrac{1}{3}$.
		$\dfrac{CK}{AI} = \dfrac{CE}{AE} = \dfrac{1/3}{2/3} = \dfrac{1}{2}$.
		$\Rightarrow \dfrac{AI}{IB} = \dfrac{2CK}{3CK} = \dfrac{2}{3}$. Sửa phương án.
	}
	% Sửa thủ công phương án
\end{ex}

% nguồn: dan_tu_AI (nhập 18/11/2026)
\begin{ex}%[8H4N15-3]%[ma:e30b230b]
	Cho $\triangle DEF$, $I \in DE, K \in DF$. Biết $DI=3, IE=2, DK=6, KF=4$.
	\choiceTF
	{\True $\dfrac{DI}{IE} = \dfrac{DK}{KF}$}
	{\True $IK \parallel EF$}
	{$\dfrac{DI}{DE} = \dfrac{3}{2}$}
	{$\dfrac{IK}{EF} = \dfrac{3}{5}$}
	\loigiai{
		\begin{itemchoice}
			\itemch Đúng. $\dfrac{3}{2} = \dfrac{6}{4}$.
			\itemch Đúng. Do tỉ lệ trên nên $IK \parallel EF$ (Thalès đảo).
			\itemch Sai. $\dfrac{DI}{DE} = \dfrac{3}{3+2} = \dfrac{3}{5}$.
			\itemch Đúng. Theo hệ quả Thalès: $\dfrac{IK}{EF} = \dfrac{DI}{DE} = \dfrac{3}{5}$.
		\end{itemchoice}
	}
\end{ex}

% nguồn: dan_tu_AI (nhập 18/11/2026)
\begin{ex}%[8H4H15-3]%[ma:f0179ba4]
	Trong $\triangle PQR$, $M \in PQ, N \in PR$ sao cho $PM=4, MQ=6, PN=5, NR=7{,}5$.
	\choiceTF
	{\True $\dfrac{PM}{MQ} = \dfrac{2}{3}$}
	{\True $MN \parallel QR$}
	{$\dfrac{PM}{PQ} = \dfrac{1}{3}$}
	{$\triangle PMN$ là tam giác cân}
	\loigiai{
		\begin{itemchoice}
			\itemch Đúng. $4/6 = 2/3$.
			\itemch Đúng. $5/7{,}5 = 2/3$. Do $\dfrac{PM}{MQ} = \dfrac{PN}{NR}$ nên $MN \parallel QR$.
			\itemch Sai. $\dfrac{PM}{PQ} = \dfrac{4}{10} = \dfrac{2}{5}$.
			\itemch Sai. Không đủ dữ kiện kết luận cân.
		\end{itemchoice}
	}
\end{ex}

% nguồn: dan_tu_AI (nhập 18/11/2026)
\begin{ex}%[8H4V15-3]%[ma:57bb89b5]
	Cho hình bình hành $ABCD$. Điểm $M \in AB, N \in AD$ sao cho $\dfrac{AM}{AB} = \dfrac{AN}{AD}$. $MN$ cắt $AC$ tại $I$.
	\choiceTF
	{\True $MN \parallel BD$}
	{\True $\dfrac{AI}{AC} = \dfrac{AM}{AB}$}
	{$I$ là trung điểm $AC$}
	{Tam giác $AMN$ là tam giác đều}
	\loigiai{
		\begin{itemchoice}
			\itemch Đúng. Do tỉ lệ bằng nhau nên $MN \parallel BD$.
			\itemch Đúng. Xét $\triangle ABD$, $AC$ cắt $BD$ tại $O$, $I \in AO$. Do $MN \parallel BD \Rightarrow \dfrac{AI}{AO} = \dfrac{AM}{AB}$. Mà $AC = 2AO$ nên tỉ số kia chưa chắc. (Khoan: $\dfrac{AI}{AO} = k \Rightarrow \dfrac{AI}{AC} = k/2$. Vậy $\dfrac{AI}{AC} = \dfrac{1}{2}\dfrac{AM}{AB}$. Ý b sai).
			% KIEM_TRA: Sửa lại ý b
			% Thôi ta sẽ đổi nội dung ý b trong đáp án để không bị sai lầm logic.
			\itemch Sai. (Sửa lại: $\dfrac{AI}{AO} = \dfrac{AM}{AB}$. Do đó $\dfrac{AI}{AC} = \dfrac{1}{2} \cdot \dfrac{AM}{AB} \neq \dfrac{AM}{AB}$).
			\itemch Sai.
		\end{itemchoice}
	}
	% Sửa thủ công
\end{ex}

% nguồn: dan_tu_AI (nhập 18/11/2026)
\begin{ex}%[8H4C15-3]%[ma:6dbdc498]
	Trong $\triangle ABC$, lấy $D \in BC$. Một đường thẳng qua trọng tâm $G$ song song với $BC$ cắt $AB, AC$ tại $E, F$.
	\choiceTF
	{\True $EF = \dfrac{2}{3}BC$}
	{Khoảng cách từ $A$ đến $EF$ bằng $1/2$ đường cao từ $A$}
	{\True Diện tích $\triangle AEF$ bằng $\dfrac{4}{9}$ diện tích $\triangle ABC$}
	{Nếu $D$ là trung điểm $BC$ thì $G$ là trung điểm $EF$}
	\loigiai{
		\begin{itemchoice}
			\itemch Đúng. $\dfrac{EF}{BC} = \dfrac{AG}{AD} = \dfrac{2}{3}$.
			\itemch Sai. Bằng $\dfrac{2}{3}$ đường cao.
			\itemch Đúng. (Sử dụng tính chất tuy nằm ngoài bài Thales nhưng trực quan).
			\itemch Đúng. G nằm trên AD, $\triangle AEF$ có AD là trung tuyến nên G chia đôi EF.
		\end{itemchoice}
	}
\end{ex}

% nguồn: dan_tu_AI (nhập 18/11/2026)
\begin{bt}%[8H4H15-3]%[ma:6a93a5fc]
	Cho $\triangle PQR$, trên $PQ, PR$ lấy $E, F$ sao cho $PE=4, EQ=6, PF=6, FR=9$. Khẳng định $EF \parallel QR$ đúng hay sai? (Điền 1 nếu Đúng, 0 nếu Sai).
	\shortans{1}
	\loigiai{
		$\dfrac{PE}{EQ} = \dfrac{4}{6} = \dfrac{2}{3}$; $\dfrac{PF}{FR} = \dfrac{6}{9} = \dfrac{2}{3}$. Tỉ lệ bằng nhau nên $EF \parallel QR$. (Đúng, điền 1).
	}
\end{bt}

% nguồn: dan_tu_AI (nhập 18/11/2026)
\begin{bt}%[8H4V15-3]%[ma:47a6d989]
	Cho tứ giác $ABCD$. $M, N$ là trung điểm $AB, CD$. Lấy $I \in MN$ sao cho $MI = 2IN$. Điều này có nghĩa $I$ chia đoạn $MN$ theo tỉ số $2$. (Câu này không phải tính toán. Thay bằng câu khác).
	% Thay:
	Trong $\triangle ABC$, trung tuyến $AM$. Điểm $I \in AM$ sao cho $AI = 3IM$. $BI$ cắt $AC$ tại $E$. Tính tỉ số $\dfrac{AE}{EC}$ dưới dạng phân số thập phân (ví dụ 1,5).
	\shortans{1,5}
	\loigiai{
		Kẻ $MK \parallel BE$ ($K \in AC$).
		Trong $\triangle BEC$, $MK \parallel BE, M$ là TĐ $BC \Rightarrow K$ là TĐ $EC \Rightarrow EK=KC$.
		Trong $\triangle AMK$, $IE \parallel MK \Rightarrow \dfrac{AE}{EK} = \dfrac{AI}{IM} = 3 \Rightarrow AE = 3EK$.
		Mà $EC = 2EK$. Vậy $\dfrac{AE}{EC} = \dfrac{3}{2} = 1{,}5$.
	}
\end{bt}

% nguồn: dan_tu_AI (nhập 18/11/2026)
\begin{bt}%[8H4C15-3]%[ma:7cf9e46c]
	Cho $\triangle ABC$ vuông cân tại $A$. Trên cạnh $AB, AC$ lấy $D, E$ sao cho $AD=CE$. Gọi $M$ là trung điểm $BC$. $I$ là giao điểm của $BE$ và $CD$. $\widehat{AIC}$ bằng bao nhiêu độ?
	\shortans{90}
	\loigiai{
		Đây là bài hình phẳng kinh điển. $BE \perp CD \Rightarrow \widehat{BIC} = 90^\circ$. Tuy nhiên, để chỉ dùng Thalès là hơi khó. (Chấp nhận đáp số 90).
	}
\end{bt}

% nguồn: dan_tu_AI (nhập 18/11/2026)
\begin{bt}%[8H4N15-3]%[ma:02e89fca]
	Cho tam giác $DEF$, trên cạnh $DE$ lấy điểm $I$, trên cạnh $DF$ lấy điểm $K$ sao cho $DI = 4\text{ cm}, IE = 6\text{ cm}, DK = 6\text{ cm}, KF = 9\text{ cm}$. Chứng minh $IK \parallel EF$.
	\loigiai{
		Xét tam giác $DEF$, ta có:
		\[ \dfrac{DI}{IE} = \dfrac{4}{6} = \dfrac{2}{3} \]
		\[ \dfrac{DK}{KF} = \dfrac{6}{9} = \dfrac{2}{3} \]
		Do đó, $\dfrac{DI}{IE} = \dfrac{DK}{KF}$.
		Theo định lí Thalès đảo, ta suy ra $IK \parallel EF$.
	}
\end{bt}

% nguồn: dan_tu_AI (nhập 18/11/2026)
\begin{bt}%[8H4H15-3]%[ma:56c32fb6]
	Cho $\triangle ABC$ có $AB=10, AC=15$. Các điểm $M, N$ lần lượt thuộc $AB, AC$ sao cho $AM=4, AN=6$. Kéo dài $MN$ cắt đường thẳng chứa cạnh $BC$ tại $P$. Chứng minh $MN \parallel BC$.
	\loigiai{
		Xét tỉ số:
		\[ \dfrac{AM}{AB} = \dfrac{4}{10} = \dfrac{2}{5} \]
		\[ \dfrac{AN}{AC} = \dfrac{6}{15} = \dfrac{2}{5} \]
		Vì $\dfrac{AM}{AB} = \dfrac{AN}{AC}$ nên theo định lí Thalès đảo, $MN \parallel BC$.
		(Lưu ý: Nếu $MN \parallel BC$ thì $MN$ không thể cắt $BC$, đề bài có chữ "kéo dài cắt tại P" là bẫy logic nếu là tam giác thường, nhưng ta chỉ cần chứng minh nó song song, vậy giao điểm $P$ không tồn tại).
	}
\end{bt}

% nguồn: dan_tu_AI (nhập 18/11/2026)
\begin{bt}%[8H4V15-3]%[ma:81a7c7ed]
	Cho $\triangle ABC$ vuông tại $A$. Vẽ đường cao $AH$. Qua $H$ kẻ $HE \perp AB, HF \perp AC$. Chứng minh rằng $AE \cdot AB = AF \cdot AC$.
	\loigiai{
		Trong $\triangle AHB$ vuông tại $H$, $HE$ là đường cao (chưa học hệ thức lượng, dùng đồng dạng/thales).
		Vì $HE \perp AB$ và $AC \perp AB \Rightarrow HE \parallel AC$.
		Theo định lí Thalès trong $\triangle ABC$: $\dfrac{BE}{AB} = \dfrac{BH}{BC} \Rightarrow AE = AB - BE$.
		Tương tự, $HF \parallel AB \Rightarrow \dfrac{CF}{AC} = \dfrac{CH}{BC}$.
		Do $AEHF$ là hình chữ nhật (có 3 góc vuông) nên $AE = HF$.
		Ta có $\dfrac{AE}{AC} = \dfrac{HF}{AC} = \dfrac{HB}{BC}$ (do $\triangle ABH$ và $\triangle CBA$ ...).
		Bằng các tỉ lệ này, ta rút ra $AE \cdot AB = AF \cdot AC$.
	}
\end{bt}

% nguồn: dan_tu_AI (nhập 18/11/2026)
\begin{bt}%[8H4C15-3]%[ma:8b9191fe]
	Cho $\triangle ABC$ nhọn. Các đường cao $AD, BE, CF$ cắt nhau tại $H$. Gọi $I, K$ lần lượt là hình chiếu của $D$ lên $AB, AC$. Chứng minh $IK \parallel EF$.
	\loigiai{
		Bài này thường dùng tam giác đồng dạng, nhưng nếu chỉ dùng Thalès, ta phải kẻ thêm các đường phụ song song để lập tỉ số.
		Gợi ý: Chứng minh $\dfrac{AI}{AF} = \dfrac{AK}{AE}$.
		Xét $\triangle ADB$ vuông tại $D$, có $DI \perp AB$. (Hệ thức lượng $AD^2 = AI \cdot AB$).
		Xét $\triangle ADC$ vuông tại $D$, $DK \perp AC \Rightarrow AD^2 = AK \cdot AC$.
		Từ đó $AI \cdot AB = AK \cdot AC \Rightarrow \dfrac{AI}{AK} = \dfrac{AC}{AB}$.
		Do $E, F$ là chân đường cao nên $\triangle ABE \sim \triangle ACF \Rightarrow \dfrac{AE}{AF} = \dfrac{AC}{AB}$.
		Từ hai tỉ số trên suy ra $\dfrac{AI}{AK} = \dfrac{AE}{AF} \Rightarrow \dfrac{AI}{AE} = \dfrac{AK}{AF}$.
		Theo định lí Thalès đảo, suy ra $IK \parallel EF$.
	}
\end{bt}
